\chapter{Abbildungen mit metrischen Voraussetzungen}
\label{b47guide:ch:maps}

In diesem Kapitel sind $(X,d_X)$ und $(Y,d_Y)$ metrische Räume.
Die Voraussetzungen an Abbildungen beziehen sich jeweils auf ihre Metriken.

\section{Stetigkeit und quantitative Abstandskontrolle}

\begin{definition}[Metrische Stetigkeit]
\label{b47guide:def:continuity}
Eine Abbildung $f:X\to Y$ heißt \emph{stetig in $x\in X$}, wenn
\[
 \forall\varepsilon>0\ \exists\delta>0\ \forall z\in X:
 \quad d_X(z,x)<\delta\ \Longrightarrow\ d_Y(f(z),f(x))<\varepsilon.
\]
Sie heißt \emph{stetig}, wenn sie in jedem Punkt von $X$ stetig ist.
\end{definition}

\begin{theorem}[Folgenkriterium der Stetigkeit]
\label{b47guide:thm:sequence-continuity}
Für $f:X\to Y$ und $x\in X$ sind äquivalent:
\begin{enumerate}
\item $f$ ist in $x$ stetig.
\item Für jede Folge $(x_n)$ in $X$ gilt
      $x_n\to x\Longrightarrow f(x_n)\to f(x)$.
\end{enumerate}
Die Grenzwerte werden bezüglich der jeweils angegebenen Metrik gebildet.
\cite{metric-notes}
\end{theorem}

\begin{definition}[Gleichmäßige Stetigkeit]
\label{b47guide:def:uniform-continuity}
Eine Abbildung $f:X\to Y$ heißt \emph{gleichmäßig stetig}, wenn
\[
 \forall\varepsilon>0\ \exists\delta>0\ \forall x,z\in X:
 \quad d_X(x,z)<\delta\ \Longrightarrow\ d_Y(f(x),f(z))<\varepsilon.
\]
Hier darf $\delta$ von $\varepsilon$, aber nicht von $x$ oder $z$ abhängen.
\end{definition}

\begin{theorem}[Folgenpaarkriterium für gleichmäßige Stetigkeit]
\label{b47guide:thm:uniform-sequences}
Eine Abbildung $f:X\to Y$ ist genau dann gleichmäßig stetig, wenn für
alle Folgen $(x_n)$ und $(z_n)$ in $X$ gilt
\[
 d_X(x_n,z_n)\to0\quad\Longrightarrow\quad
 d_Y(f(x_n),f(z_n))\to0.
\]
Die beiden Folgen selbst brauchen dabei keinen Grenzwert zu besitzen.
\end{theorem}
\begin{definition}[Lipschitz- und bi-Lipschitz-Abbildung]
\label{b47guide:def:lipschitz}
Eine Abbildung $f:X\to Y$ heißt \emph{$L$-Lipschitz} für eine reelle
Zahl $L\geq0$, wenn
\[
 d_Y(f(x),f(z))\leq Ld_X(x,z)\qquad(x,z\in X).
\]
Sie heißt \emph{bi-Lipschitz auf ihr Bild}, wenn reelle Konstanten
$0<c\leq C<\infty$ existieren mit
\[
 c\,d_X(x,z)\leq d_Y(f(x),f(z))\leq C\,d_X(x,z)
 \qquad(x,z\in X).
\]
\end{definition}

\begin{theorem}[Stufen der Abstandskontrolle]
\label{b47guide:thm:continuity-hierarchy}
Jede Lipschitz-Abbildung ist gleichmäßig stetig; jede gleichmäßig
stetige Abbildung ist stetig. Die Umkehrungen gelten im Allgemeinen
nicht. Die Abbildung $x\mapsto\sqrt{x}$ auf $[0,1]$ mit dem üblichen
Abstand ist gleichmäßig stetig, aber nicht Lipschitz.
Die Abbildung $x\mapsto x^2$ auf $\mathbb R$ mit dem üblichen Abstand
ist stetig, aber nicht gleichmäßig stetig.
\end{theorem}

\begin{definition}[Isometrische Einbettung und Isometrie]
\label{b47guide:def:isometry}
Eine Abbildung $i:X\to Y$ heißt \emph{isometrische Einbettung}, wenn
\[
 d_Y(i(x),i(z))=d_X(x,z)\qquad(x,z\in X).
\]
Eine surjektive isometrische Einbettung heißt hier \emph{Isometrie
zwischen $X$ und $Y$}. Der Hauptband verwendet auch für eine
isometrische Einbettung die Bezeichnung \emph{isometrische Abbildung}.
\end{definition}

\begin{theorem}[Metrische Invarianten]
\label{b47guide:thm:metric-invariants}
Eine isometrische Einbettung ist injektiv und besitzt auf ihrem
Bild eine isometrische Umkehrabbildung. Sie erhält und reflektiert
Konvergenz, die Cauchy-Eigenschaft, Beschränktheit, totale
Beschränktheit, Vollständigkeit und Kompaktheit, jeweils zwischen
$X$ und dem metrischen Teilraum $i[X]$.
Auch bi-Lipschitz-Abbildungen auf ihr Bild erhalten und reflektieren
die genannten Eigenschaften. Ihre Umkehrabbildung auf dem Bild ist
bi-Lipschitz.
Insbesondere ist das isometrische Bild eines vollständigen Raums
in jedem metrischen Zielraum abgeschlossen.
\end{theorem}

\section{Erhaltung und Kompaktheit}

\begin{theorem}[Gleichmäßige Stetigkeit und Cauchy-Struktur]
\label{b47guide:thm:uniform-preservation}
Ist $f:X\to Y$ gleichmäßig stetig, so bildet $f$ jede Cauchy-Folge
in $X$ auf eine Cauchy-Folge in $Y$ ab. Für jede total beschränkte
Teilmenge $A\subseteq X$ ist $f[A]$ total beschränkt.
Vollständigkeit des Bildes folgt hieraus im Allgemeinen nicht.
\cite{metric-notes}
\end{theorem}

\begin{theorem}[Stetige Bilder metrisch kompakter Mengen]
\label{b47guide:thm:compact-image}
Ist $f:X\to Y$ stetig und $K\subseteq X$ kompakt, so ist $f[K]$
kompakt. Insbesondere ist $f[K]$ vollständig, total beschränkt
und in $Y$ abgeschlossen.
\cite{metric-notes}
\end{theorem}

\begin{theorem}[Heine--Cantor]
\label{b47guide:thm:heine-cantor}
Ist $X$ kompakt, so ist jede stetige Abbildung $f:X\to Y$
gleichmäßig stetig. Eine zusätzliche Vollständigkeits- oder
Kompaktheitsannahme an $Y$ ist hierfür nicht erforderlich.
\cite{metric-notes}
\end{theorem}

